\documentclass[UTF8]{ctexart}

\usepackage{amscd,amsthm,amsfonts,amssymb,amsmath}
\usepackage[colorlinks=true]{hyperref}
\usepackage[all]{xy}
\usepackage{mathrsfs}
\usepackage{tikz}
\usepackage{bm}
\usepackage{geometry}
\usepackage{xcolor}
\usepackage{eso-pic}
% \usepackage{CJK}
\usepackage{arydshln}
% \usepackage[11pt]{moresize}



\newtheorem{thm}{定理}
\newtheorem{cor}[thm]{推论}
\newtheorem{lem}[thm]{引理}
\newtheorem{prop}[thm]{命题}
\newtheorem{defn}[thm]{定义}
\newtheorem{ques}[thm]{问题}
\newtheorem{eg}[thm]{例题}
\newtheorem{ex}[thm]{习题}
\newtheorem{rmk}[thm]{注释}
\newtheorem{ntn}[thm]{记号}




\begin{document}


\title{习题课材料（七）}
\date{}

\maketitle


%\noindent{\Large\textbf{注: 带 $\heartsuit $号的习题有一定的难度、比较耗时, 请量力为之.}}

\begin{ex}
  奇数阶反对称矩阵不可逆.
\end{ex}

\begin{proof}
  事实上, \(\det A = \det (-A^{\mathrm{T}}) = \det(-A) = (-1)^{n}\det A\).
  故若 \(\det A \neq 0\), 必有 \((-1)^{n} = 1\), 即 \(n\) 为偶数.
\end{proof}

\begin{ex}
  对分块上三角矩阵 \(X=\begin{bmatrix}A & C \\ O & B\end{bmatrix}\),
  证明 \(\det X = \det A \cdot \det B\).
\end{ex}

\begin{proof}
  若 \(\det A = 0\),
  则 \(A\) 的列向量线性相关. 因此 \(X\) 的列线性相关, 从而
  \(\det X = 0 = \det A\det B\).

  若 \(\det A \neq 0\),
  则 \(A\) 可以写成初等矩阵的乘积
  \(A = E_1\cdots E_{m}\).
  利用分块矩阵乘法, 可知
  \begin{equation*}
    X =
    \begin{bmatrix}
      E_1  \\ & I
    \end{bmatrix} \cdots
    \begin{bmatrix}
      E_m  \\ & I
    \end{bmatrix}
    \cdot
    \begin{bmatrix}
      I & C' \\ & B
    \end{bmatrix},
  \end{equation*}
  其中\(C' = E_{m}^{-1}\cdots E_{1}^{-1}C\).
  由于 \(\det A = \prod_{i=1}^{m} \det E_i\),
  因此只需证明 \(\det\begin{bmatrix}I& C'\\ & B\end{bmatrix}=\det B\).
  但是此矩阵为 \(\mathrm{diag}(I,B)\) 乘以一系列第三类初等矩阵, 故而行列式与
  \(\mathrm{diag}(I,B)\) 相同, 等于 \(\det B\).
\end{proof}

\begin{ex}
  设 \(A\), \(B\) 分别为 \(m\times n\), \(n\times m\) 矩阵.
  证明 \(\det(I_m + AB) = \det(I_n+BA)\).
\end{ex}

\begin{proof}
  考虑分块矩阵
  \begin{equation*}
    X=
    \begin{bmatrix}
      I_m & -A \\ B & I_n
    \end{bmatrix}.
  \end{equation*}
  注意到
  \begin{equation*}
    \begin{bmatrix}
      I_m   \\ -B & I_n
    \end{bmatrix}
    \cdot X
    =
    \begin{bmatrix}
      I_m & -A \\ & I_n + BA
    \end{bmatrix},
    \quad
    \begin{bmatrix}
      I_m & A\\   & I_n
    \end{bmatrix}
    \cdot X
    =
    \begin{bmatrix}
      I_m+AB \\ B & I_n
    \end{bmatrix}.
  \end{equation*}
  因此 \(\det(I_m+AB)=\det X=\det(I_n+BA)\).
\end{proof}

\begin{ex}[\(\heartsuit\)]
  设 \(T = [\boldsymbol{t}_1,\ldots,\boldsymbol{t}_n]\) 为 \(n\) 阶方阵.
  证明 Hadamard 不等式
  \[
    |\det T| \leq \|\boldsymbol{t}_1\| \cdots \|\boldsymbol{t}_n\|.
  \]
\end{ex}

\begin{proof}
  无妨设 \(\boldsymbol{t}_1,\ldots\boldsymbol{t}_n\) 线性无关, 否则不等式显然成立.
  令 \(T = QR\) 为 \(T\) 之 QR 分解, 其中
  \(Q = [\boldsymbol{q}_1,\ldots,\boldsymbol{q}_n]\) 为正交方阵,
  \(R\) 为上三角方阵. 令 \(R = (a_{ij})\),
  则 \(\det R = \prod_{i=1}^{n} a_{ii}\).

  将矩阵等式 \(T=QR\)  写作向量之间的等式:
  \begin{equation*}
    \boldsymbol{t}_{j} = \sum_{i\leq j} a_{ij} \cdot \boldsymbol{q}_i.
  \end{equation*}
  由于 \(\boldsymbol{q}_1, \boldsymbol{q}_2, \ldots,\boldsymbol{q}_n\)
  构成 \(\mathbb{R}^n\) 之单位正交基,
  由勾股定理, 可知 \(\|\boldsymbol{t}_j\|^{2} = \sum |a_{ij}|^2 \geq |a_{jj}|^{2}\).
  又 \(|\det Q| = 1\), 故
  \begin{equation*}
    |\det T| = |\det R| = \prod_{i=1}^{n} |a_{ii}| \leq \prod_{i=1}^{n}\|t_i\|. \qedhere
  \end{equation*}
\end{proof}


\begin{ex}
  计算行列式
  \begin{equation*}
    \begin{vmatrix}
      1 & 0 & 0 & 2 \\
      0 & 3 & 4 & 5 \\
      5 & 4 & 0 & 3 \\
      2 & 0 & 0 & 1
    \end{vmatrix}
  \end{equation*}
\end{ex}

\begin{proof}[解]

  第一行的 $0$ 较多, 故按第一行展开.
  所求之数为
  \begin{equation*}
    1 \cdot
    \begin{vmatrix}
      3 & 4 & 5 \\
      4 & 0 & 3 \\
      0 & 0 & 1
    \end{vmatrix}
    - 2 \cdot
    \begin{vmatrix}
      0 & 3 & 4 \\
      5 & 4 & 0 \\
      2 & 0 & 0
    \end{vmatrix}.
  \end{equation*}
  对第一个 3 阶行列式, 按第三行展开:
  \begin{equation*}
    \begin{vmatrix}
      3 & 4 & 5 \\
      4 & 0 & 3 \\
      0 & 0 & 1
    \end{vmatrix}
    = (-1)^{3+3} \cdot
    \begin{vmatrix}
      3 & 4 \\ 4 & 0
    \end{vmatrix}
    = -16.
  \end{equation*}
  对第二个 3 阶行列式, 也按第三行展开:
  \begin{equation*}
    \begin{vmatrix}
      0 & 3 & 4 \\
      5 & 4 & 0 \\
      2 & 0 & 0
    \end{vmatrix}
    =(-1)^{3+1} \cdot 2 \cdot
    \begin{vmatrix}
      3 & 4 \\ 4 & 0
    \end{vmatrix} = - 2\cdot 16.
  \end{equation*}
  因此, 最后答案为 \((-16)-4(-16)=48\).
\end{proof}

\begin{ex}
  \(\begin{vmatrix} 2x & x & 1 & 2 \\ 1 & x & 1 & -1\\ 3 & 2 & x & 1 \\ 1 & 1 & 1& x\end{vmatrix}\)
  中 \(x^4\), \(x^3\) 系数.
\end{ex}

\begin{proof}[答案]
  \(x^4\) 系数为 \(2\), \(x^3\) 系数为$ -1$. 可考察行列式的完全展开式,
  也可按照第一行展开, 然后再利用三阶行列式的展开式逐一计算.
\end{proof}

\begin{ex}
  设方阵 \(A,B\) 满足 \(AB=BA\). 证明
  \(\det\begin{bmatrix}A & B\\-B & A\end{bmatrix}=\det(A^2+B^2)\).
\end{ex}

\begin{proof}
  令 \(\mathrm{i}=\sqrt{-1}\) 为虚数单位.
  利用分块矩阵乘法直接计算可知
  \begin{equation*}
    \begin{bmatrix}
      \mathrm{i}I & I \\ I & \mathrm{i}I
    \end{bmatrix}
    \begin{bmatrix}
      A & B \\ -B & A
    \end{bmatrix}
    \begin{bmatrix}
      \mathrm{i}I & I \\ I & \mathrm{i}I
    \end{bmatrix}^{-1} =
    \begin{bmatrix}
      A+\mathrm{i}B & \\ & A-\mathrm{i}B
    \end{bmatrix}.
  \end{equation*}
  因此所求行列式等于 \(\det((A-\mathrm{i}B)(A+\mathrm{i}B))\).
  最后, 由于 \(AB=BA\),
  \(\det((A-\mathrm{i}B)(A+\mathrm{i}B))=\det(A^2+B^2)\).
\end{proof}

\begin{ex}
  \begin{enumerate}
  \item
    设 \(A\) 是正交矩阵, \(\det A < 0\). 证明 \(I_n + A\) 不可逆.
  \item   设 \(A\) 为奇数阶正交矩阵, \(\det A > 0\).
    证明 \(I_n - A\) 不可逆.
  \end{enumerate}
\end{ex}

\begin{proof}
  (1) 今 \(\det(I_n+A)\det(A^{\mathrm{T}}) = \det(A^{\mathrm{T}}+I_n)=\det(I_n+A)\).
  若 \(\det(I_n +A)\neq 0\), 则 \(\det(A^{\mathrm{T}})=\det A = 1 > 0\).

  (2)   今 \(\det(I_n -A) \det(A^{\mathrm{T}})=\det(A^{\mathrm{T}}-I_n)=\det(A-I_n)\).
  由于 \(n\) 为奇数, \(\det(A-I_n)=-\det(I_n-A)\). 若 \(\det(I_n-A)\neq 0\),
  则 \(\det(A^{\mathrm{T}})=\det(A)=-1<0\).
\end{proof}

\begin{ex}[\(\heartsuit\)]
  设 \(A\) 为对角占优方阵, 对角元素全为正数. 证明 \(\det A > 0\).
\end{ex}

\begin{proof}
  令 \(A(\lambda) = \lambda I_n + A\).
  根据行列式展开式, \(\det(A(\lambda))\) 为 \(\lambda\) 的首项系数为 1 的 \(n\) 次多项式.
  由于对任何 \(\lambda \geq 0\), \(A(\lambda)\) 对角占优,
  故对任何 \(\lambda \geq 0\), \(\det (A(\lambda)) \neq 0\).
  显然 \(\det(A(\lambda)) \to +\infty\), (\(\lambda\to +\infty\)).
  又由于 \(\det(A(\lambda))\) 连续, 在 \(\{0\leq \lambda < +\infty\}\)
  上处处非零, 故介值定理推出, 对任何 \(\lambda \geq 0\),
  \(\det(A(\lambda))\geq 0\). 由于 \(\det(A(0))=\det(A)\neq 0\),
  必有 \(\det A > 0\).
\end{proof}

\begin{ex}
  设 \(A\) 是元素为整数的可逆矩阵. 证明, 若 \(A\) 可逆, 并且 \(A^{-1}\)
  的元素也为整数的必要且充分条件是 \(|\det A| =  1\).
\end{ex}

\begin{proof}

  若 \(\det A = \pm 1\),
  令 \(\mathrm{adj}(A)\) 为 \(A\) 的伴随矩阵,
  即 \(\mathrm{adj}(A)^{\mathrm{T}}\) 的 \((i,j)\) 元素为 \(a_{ij}\)
  的代数余子式.
  注意到 \(\mathrm{adj}(A)\) 的元素都是整数.
  则 \(A \cdot \mathrm{adj}(A) = \pm I\), 因此, \(\pm\mathrm{adj}(A)\)
  为 \(A\) 的逆, 其元素俱为整数.

  反之, 若 \(AC = \mathrm{I}\), \(C\) 之元素为整数. 则 \(\det A\), \(\det C\)
  都是整数, 且
  \(\det A \cdot \det C = 1\). 故而必有 \(|\det A| = 1\).
\end{proof}
\end{document}
